The verification files distinguish four checks.
\begin{enumerate}
\item \textbf{Semantic interface binding.} An independent literal audit checks the doubled frame, every actual seed and its principal norm, the complete private anchor, the merged aggregate row, all six power requests, all eleven modular requests, and the shared factorial/binomial value. It verifies that every needed exponent is in the semantic frame and that no private Pell witness enters that frame.
\item \textbf{Signed-domain recovery.} The signed certificate records the initial scalar predicates, every polynomial sign implication, the exact twelve product predicates, the thirty recovered root coordinates, all deferred quantities, and the noncircular order in which their signs become available. Each guard is checked as an actual row with fresh private witnesses.
\item \textbf{Complete equation transformations.} The natural elimination certificates and the integer monic inverse certificates are replayed against whole residual lists. The final coordinate translation is verified exactly. An independent expansion of the final residuals is compared coefficient by coefficient with the complete collected polynomial.
\item \textbf{Ordinary theorem diagnostics.} The separate small Pell program checks signed and zero-index branches, constructive witnesses, and an explicit negative-quotient counterexample. Those finite checks supplement the proof in Section~\ref{sec:one-row-pell}; they do not certify its unbounded quantifiers by enumeration.
\end{enumerate}

The final receipts state the exact scopes executed, artifact hashes, and counts. Normal and optimized Python are both supported; acceptance-critical checks use explicit exceptions. The portable reconstruction is checked in a fresh output directory against the retained equation hashes. There is no search for a ZFC refutation or a zero of the final polynomial in these verification programs.
